\par\clearpage
\noindent\makebox[0pt][l]{\fontsize{13}{16}\selectfont \textbf{382}}\hfill{\small\textbf{BRIEFWECHSEL MIT GERLING.}}\hfill\null\par\medskip
\begin{center}
\textsc{Gauss} an \textsc{Gerling}.\quad Göttingen, 27.~Juli 1823.
\end{center}

\par\noindent\hspace*{\parindent}\mbox{.\ .\ .\ .\ .\ .}\quad Ich habe mir von \textsc{Repsold} noch ein neues Okular für meinen Theodoliten anfertigen lassen, mit Spinnenfäden von fast unglaublicher Feinheit, die \(29^{\prime\prime}\) voneinander entfernt sind.
(Die Fäden im Okular von \textsc{Ertel}, die auch schon recht fein, aber viel dicker als jene sind, sind \(36^{\prime\prime}\) voneinander entfernt.)
Bei der Beobachtung entfernter \textls[120]{blasser} irdischer Gegenstände ist diese Einrichtung äußerst vorteilhaft. Leider bekam ich sie erst post festum.
Der Hamburger Turm, der wegen der vom Höhenrauch getrübten Luft von Nindorf, Timpenberg und Lüneburg aus fast immer sehr blass aussah, hat mir sehr zu schaffen gemacht, ebenso der Lüneburger Turm von Hamburg aus. Im Allgemeinen sind die Messungen in diesen Gegenden nicht ganz so präzise wie die früheren, wozu übrigens auch das ständige Schwanken der Türme beigetragen hat.
Der Michaelis-Thurm in Hamburg war, solange ich dort war, nie ruhig. Seine horizontalen, pendelartigen Schwankungen gehen oft über \(\frac{1}{2}\) Bogenminute hinaus.
Dieser Turm war von allen meinen Dreieckspunkten als Standpunkt und (vom Brocken abgesehen) auch als Zielpunkt mit Abstand der schlechteste.
\textsc{Repsold} fertigt auch neue Libellen von außerordentlicher Vollkommenheit an, die er statt mit Weingeist mit Naphtha füllt.
Sie kommen sehr viel schneller zur Ruhe als die anderen.
